\documentclass{article}
\input{../../../lib.sty}

\title{\texttt{fn make\_vec}}
\author{Michael Shoemate}
\date{}

\begin{document}

\maketitle

\contrib
Proves soundness of the implementation of \rustdoc{transformations/fn}{make\_vec} in \asOfCommit{mod.rs}{f5bb719}.

This transformation simply wraps an input scalar in a singleton vec.
The output metric then becomes an Lp distance.
\section{Hoare Triple}
\subsection*{Precondition}
\subsubsection*{Compiler-Verified}
\begin{itemize}
    \item Generic \texttt{T} implements trait \rustdoc{traits/trait}{Number}
    \item Generic \texttt{Q} implements trait \rustdoc{traits/trait}{Number}
\end{itemize}

\subsubsection*{User-Verified}
None

\subsection*{Pseudocode}
\lstinputlisting[language=Python,firstline=2,escapechar=|]{./pseudocode/make_vec.py}

\subsection*{Postcondition}
\begin{theorem}
    \validTransformation{(\texttt{input\_space, T, Q})}{\texttt{make\_vec}}
\end{theorem}

\begin{proof}
    The function is infallible, and the output domain trivially follows, since all output vectors are length-one.
    The function is 1-stable because:

    
    \begin{align}
        \max_{x \sim x'} d_{Lp}(f(x), f(x'))\\
        = \max_{x \sim x'} (\sum_i (x_i - x_i')^p)^{1/p}\\
        = \max_{x \sim x'} ((x_1 - x_1')^p)^{1/p}\\
        = \max_{x \sim x'} |x_1 - x_1'|\\
        = \max_{x \sim x'} d_{Abs}(x_1, x_1')\\
        = 1 \cdot \texttt{d\_in}
    \end{align}
\end{proof}

\end{document}
